\subsection{Algebraic interpolation with linear logarithmic section size}

\begin{theorem}[DP60: bounded representatives and section norms]\label{DP60}
Let $E=\mathbf Q(\zeta_7)$, let $\mathfrak p\mid11$, and put
\[
 A_0=\mathcal O_E[1/14],\qquad
 R_0=\mathcal O_{E,\mathfrak p},\qquad
 R=\widehat R_0=W(\mathbf F_{11^3}).
\]
Use the good Jacobian model $J/A_0$, the trace-one ample line $\mathcal L$, and the invariant product normalization of DP49 and DP54:
\[
 \Theta_s=f_sU^\kappa B^{-\kappa},
 \qquad U\in1+(z)R[[z]],\qquad B=\exp b.
\]
Fix continuous Hermitian metrics on $\mathcal L$ at the complex embeddings and use their tensor powers.
There are constants $a\ge1$, $b_0\ge0$, $C_1,C_2,C_3$, independent of $n$ and $\kappa$, such that for every $n\ge1$ and every $\kappa\ge an+b_0$ with $11\nmid\kappa$, there is a section
\[
 0\ne s_{n,\kappa}\in H^0(J_{A_0},\mathcal L^{\otimes\kappa})
\]
satisfying
\begin{align}
 f_{s_{n,\kappa}}U^\kappa
 &\equiv1\pmod{11R[[z]]+(z)^{n+1}},
 \label{eq:boundedrepresentativejet}\\
 \log\max\{1,\|s_{n,\kappa}\|_{\sigma,\sup}\}
 &\le C_1\kappa+C_2\log(n+1)+C_3
 \label{eq:linearsectionsize}
\end{align}
for every complex embedding $\sigma$.
Every explicit ordinary witness of DP58 of degree at most $n$ retains its exact pure-kernel valuation on this section.
The same holds for the explicit first-row two-column harmonic witnesses, using two copies of the section.

In particular one can choose
\[
 \kappa=M^2=an+O(\sqrt n),\qquad 11\nmid M,
\]
so both the tensor exponent and the logarithmic section sup-size are $O(n)$, while the exact selected valuations, including the unit factor $M^{-m}$ at each selected degree $m\le n$, are preserved.
The sections are integral at every finite place away from $2$ and $7$ in this fixed model.
For fixed tensor-compatible adelic metrics induced by the model outside these bad places, their adelic logarithmic sup-size is also $O(\kappa+\log(n+1))$.
\end{theorem}

\begin{proof}
\emph{Constant precision preserves the explicit jets.}
First work at any odd $p$ with the actual pair of DP46.
Write $s_*=1/(p-1)$ and, for a primitive line of intrinsic length $k$, put
\[
 N=p^k,\qquad w=\frac{s_*N}{N+1},\qquad
 c=\frac{s_*}{N+1}.
\]
At $m=(N+1)p^r$, the pure unit-parameter jet has valuation $cm-s_*$ by DP58.
If an integral prefactor is congruent to one modulo $p$ and order $n+1$, write
\[
 W=1+pV+H,\qquad V\in R[[z]],\quad H\in(z)^{n+1}.
\]
For $m\le n$, the term $H$ makes no contribution after inverse-logarithm substitution, since that substitution preserves total order.
DP46 makes the substituted integral function $V$ times the unit-parameter kernel integral at weight $w$.
As $m$ has digit sum two, the degree-$m$ error derivative therefore has valuation at least
\[
 1-wm+v_p(m!)=1+cm-2s_*>cm-s_*.
\]
The inequality follows from $s_*\le1/2$.
Thus reduction modulo $p$ is sufficient; the exact value is preserved for every unit exponent, including $-\kappa$.

For two copies, modulo total order $n+1$ one has
\[
 (1+pV_1)(1+pV_2)-1
 =p(V_1+V_2+pV_1V_2).
\]
The bracket is an integral function on the product formal group.
The full-polydisc argument of DP51 applies to any such function: the contracted inverse-logarithm coordinates and the product exponential are both integral.
At the first-row critical common weight
$w_*=s_*p^{h-1}/(p^{h-1}+1)$, with $c_*=s_*-w_*$, take
\[
 a_r=p^r,\qquad b_r=p^{h-1+r},\qquad m=a_r+b_r.
\]
For every supported differential monomial of $P_{a,b}$, its coefficient times its multi-factorial is, up to sign,
\[
 \binom ai\binom bi(a-i)!(b-i)!(i!)^2=a!b!.
\]
Consequently the full harmonic error has valuation at least
\[
 1-w_*m+v_p(a_r!b_r!)
 =1+c_*m-2s_*>c_*m-s_*.
\]
This is strictly above the exact pure value in DP58.
All product errors are covered, including nonseparable prefactors: the difference of products is $p$ times an integral product-coordinate function.

\emph{A finite set of global generators.}
Return to $p=11$.
Over $R_0$, choose the projective-coordinate sections from DP56,
\[
 h_0,h_1,h_2,h_3\in H^0(J_{R_0},\mathcal L^{\otimes a}),
\]
with $h_0(e)$ a unit, $h_i(e)=0$ for $i>0$, and $t_i=h_i/h_0$ integral formal parameters.
Choose $b$ so every power $\mathcal L^{\otimes m}$ with $m\ge b$ is globally generated.
For each residue class $\ell\bmod a$, choose one exponent $b_\ell\ge b$ in that class and one section
\[
 \eta_\ell\in H^0(J_{R_0},\mathcal L^{\otimes b_\ell})
\]
whose value at zero is a unit.
These choices use the same known ample-power and global-generation statements as DP56, namely Stacks Project Tags 01VT and 01Q3.
Put $b_0=\max_\ell b_\ell$.
For $\kappa\ge an+b_0$, choose $\ell\equiv\kappa\bmod a$ and set
\[
 q=\frac{\kappa-b_\ell}{a}\ge n.
\]
The section $\eta_\ell h_0^q$ is an integral local frame at zero.
Only finitely many auxiliary generators have been chosen.

Flat cohomology base change for $A_0\to R_0$ identifies these section modules over $R_0$ with the localizations of their modules over $A_0$.
The precise known statement is Stacks Project Lemma~30.5.2(2), Tag~02KH.\footnote{\url{https://stacks.math.columbia.edu/tag/02KH}.}
Clear the finitely many denominators using algebraic integers outside $\mathfrak p$.
One may do this integrally because a denominator in $A_0\setminus\mathfrak p$ has a numerator in $\mathcal O_E\setminus\mathfrak p$, and only powers of $14$ were inverted.
Use a common clearing scalar for all $h_i$, preserving their ratios.
The scalars are units over $R_0$, so the unit values, vanishing and formal-coordinate properties persist.
The resulting fixed generators are global sections over $A_0$; clearing can add divisibility at other primes but introduces no new denominator.

\emph{Bounded integral coefficient representatives.}
For every $C\ge1$ the map
\[
 \mathcal O_E/p^C\mathcal O_E\longrightarrow R/p^CR
\]
is surjective.
Indeed $p\mathcal O_E=\mathfrak p\overline{\mathfrak p}$, and the Chinese remainder theorem identifies the source with
\[
 \mathcal O_E/\mathfrak p^C\times
 \mathcal O_E/\overline{\mathfrak p}^{\,C}.
\]
Localization does not change the first quotient, since all denominators outside $\mathfrak p$ are units there, and completion preserves it as well.
Only $C=1$ is needed.
As $11$ has order three modulo $7$, the residue of $\zeta_7$ has degree three, and the fixed set
\[
 \mathcal S=\{c_0+c_1\zeta_7+c_2\zeta_7^2:
                 0\le c_i\le10\}
\]
represents $R/pR$ bijectively.
Every element is an algebraic integer and has absolute value at most $30$ at every complex embedding.

\emph{The actual global interpolant.}
Fix an integral formal frame $\tau$ of $\mathcal L$ and let
\[
 g_\kappa=\frac{\eta_\ell h_0^q}{\tau^\kappa}
 \in R[[z]]^\times.
\]
Express the integral unit $g_\kappa^{-1}U^{-\kappa}$ in the integral parameters $t_i$, truncate through degree $n$, and write the result as
\[
 \sum_{|\alpha|\le n}q^*_\alpha t^\alpha,
 \qquad q^*_\alpha\in R.
\]
Choose $q_\alpha\in\mathcal S$ congruent to $q^*_\alpha$ modulo $p$, and define
\begin{equation}\label{eq:boundedglobalinterpolant}
 s_{n,\kappa}
 =\eta_\ell\sum_{|\alpha|\le n}q_\alpha
      h_0^{q-|\alpha|}h_1^{\alpha_1}
      h_2^{\alpha_2}h_3^{\alpha_3}.
\end{equation}
All powers are nonnegative and each summand has tensor exponent $b_\ell+aq=\kappa$.
The generators are integral over $A_0$ and the coefficients lie in $\mathcal O_E$, so this is a global section over $A_0$.
Its formal coefficient multiplied by $U^\kappa$ is congruent to one modulo $p+(t)^{n+1}$, hence modulo $p+(z)^{n+1}$.
This proves \eqref{eq:boundedrepresentativejet} and makes the section nonzero, with unit value at the chosen prime.
The constant-precision argument proves all the selected jet claims, simultaneously for every witness of degree at most $n$.
The construction already lies over the number field; no uncontrolled final approximation in a completed section space is used.

\emph{Norms and all finite places.}
Choose $H\ge1$ bounding the sup norms of the finitely many $h_i$ and $\eta_\ell$ at all complex embeddings in their indicated tensor metrics.
There are $\binom{n+3}{3}$ terms in \eqref{eq:boundedglobalinterpolant}.
Tensor compatibility and the triangle inequality give the explicit estimate
\[
 \|s_{n,\kappa}\|_{\sigma,\sup}
 \le30\binom{n+3}{3}H^{q+1}.
\]
Since $q\le\kappa/a$, its logarithm is bounded by
\[
 \log30+\log\binom{n+3}{3}
       +(\kappa/a+1)\log H,
\]
proving \eqref{eq:linearsectionsize}.
At every model place outside $2,7$, integrality gives sup norm at most one.
For any fixed continuous metrics on the proper analytifications at the remaining finite places, bound the same finite generator set by constants $H_v\ge1$.
The coefficients are algebraic integers, so the ultrametric inequality gives
\[
 \|s_{n,\kappa}\|_{v,\sup}\le H_v^{q+1}.
\]
The normalized adelic sum of positive logarithmic sup norms is therefore $O(\kappa+\log(n+1))$.
There are no other finite-place contributions and no new good-prime denominators.

Finally choose $m_0=\lceil\sqrt{an+b_0}\rceil$ and let $M=m_0$ unless $11\mid m_0$, in which case set $M=m_0+1$.
Then $11\nmid M$ and $M^2=an+O(\sqrt n)$.
The preceding construction works at $\kappa=M^2$, and its normalized kernel exponent $-\kappa$ is a unit.
DP58 is uniform over all such units.
The factor $M^{-m}$ has valuation zero at every selected derivative degree.
This proves the linear bounds and exact-value assertions.
\end{proof}

Two fixed continuous Hermitian metrics on the line over a compact complex fibre differ by bounded positive factors; on its $\kappa$th tensor power the logarithmic discrepancy is $O(\kappa)$.
Thus the theorem also applies to any fixed geometric canonical theta metric.
In a fixed honest holomorphic frame on a relatively compact chart, dividing by the frame and multiplying by a fixed holomorphic gauge $\exp(\kappa h_L)$ costs only another $\exp(O(\kappa))$ factor.
This statement does not assign an archimedean norm to the $p$-adic display frame or identify it with a source complex gauge.
Level-dependent normalizations require their own comparison.

The result is a construction in the full algebraic section space, with exact finite local values and upper bounds for fixed norms.
It does not identify the prescribed torsion-translate span, force a prescribed degree sequence to lie on the explicit witness ray, or compare the constants with a global gain-minus-cost budget.
No conclusion about existence or impossibility of an entire shorter lift chain follows from this norm bound alone.
